% Generated draft: requires mathematical and source review.
\documentclass[11pt]{article}
\usepackage[margin=25mm]{geometry}
\usepackage{fontspec}
\setmainfont{Latin Modern Roman}
\usepackage{amsmath,amssymb,mathtools,unicode-math}
\setmathfont{Latin Modern Math}
\usepackage{enumitem,needspace,xurl,hyperref}
\setlength{\parindent}{0pt}
\setlength{\parskip}{5pt}
\setlength{\emergencystretch}{3em}
\newcommand{\N}{\mathbb{N}}
\newcommand{\Z}{\mathbb{Z}}
\newcommand{\Q}{\mathbb{Q}}
\newcommand{\R}{\mathbb{R}}
\newcommand{\C}{\mathbb{C}}
\newcommand{\F}{\mathbb{F}}
\newcommand{\A}{\mathbb{A}}
\newcommand{\PP}{\mathbb{P}}
\newcommand{\E}{\mathbb{E}}
\newcommand{\eps}{\varepsilon}
\newcommand{\dd}{\,\mathrm{d}}
\begin{document}
\section*{UnsolvedMath ID: 30005880}
\textbf{Source problem code:} catalogue entry
\subsection{Definitions and mathematical statement}
Let \(F\) be a finite field, let \(G\) be a linear algebraic group defined over \(F\), and write \(H=G(F)\), \(r=\operatorname{rank}(G)\), and \(d=\dim G\). Let \(g\in H\) be regular semisimple, and let \(A\subseteq H\) satisfy \(\langle A\rangle=H\). For an integer \(k\geq 1\), define
\[
A^k=\{a_1\cdots a_k:a_1,\ldots,a_k\in A\},\qquad
\operatorname{Cl}_H(g)=\{hgh^{-1}:h\in H\},
\]
and define the centralizer \(C_H(g)=\{h\in H:hg=gh\}\). No symmetry assumption on \(A\), or assumption that \(1\in A\), is supplied.

The requested inequality, for \(d>0\), is
\[
\bigl|A\cap\operatorname{Cl}_H(g)\bigr|
\leq C\,|A^k|^{\,1-r/d},
\]
with \(C>0\) and \(k\geq 1\) as small as possible. The rank-uniform interpretation suggested by the source is to find functions \(C(r)\) and \(k(r)\) for which this holds for every admissible \(F,G,g,A\) of rank \(r\), independently of the field size and the particular generating set. In particular, the source asks whether both constants can have polynomial rank dependence: whether there exist absolute constants \(c>0\) and \(b\geq 0\) such that, for every admissible rank \(r\geq 1\), one can take
\[
C(r)\leq c r^b,\qquad 1\leq k(r)\leq c r^b.
\]
It also asks whether character theory can help establish such bounds; this is a methodological question without a separately specified mathematical conclusion.

\emph{Unresolved scope qualifications.} The source does not specify connectedness, reductivity, simplicity, or whether rank means absolute rank or \(F\)-split rank. In the customary connected reductive setting, absolute rank is the dimension of a maximal torus, and a semisimple element is regular when its algebraic centralizer has dimension equal to that rank. These conventions describe a possible precise setting; they are not additional hypotheses supplied by the record. The admissible class and conventions must therefore be confirmed before the displayed rank-uniform formulation can be treated as an unqualified problem statement. The source also does not address \(d=0\), when its exponent is undefined, or specify the intended treatment of rank zero.

The source explicitly describes a conjugacy class in \(G(F)\), but its parenthetical alternative repeats the same notation in the supplied transcription. If a geometric conjugacy class was intended there, the alternative would involve
\[
\operatorname{Cl}_{G(\overline F)}(g)
=\{xgx^{-1}:x\in G(\overline F)\},
\]
intersected with \(A\subseteq G(F)\). This alternative is recorded conditionally because the transcription does not establish it; rational and geometric conjugacy classes must not be silently identified.

\emph{Supplied context.} The report passage appears to give a currently known bound \(C,k=O(r)^{O(r^2)}\), although the supplied extraction has damaged superscripts. It says that using the conjugacy-class structure improves the rank dependence compared with available bounds for intersections with arbitrary varieties, referring to an unidentified reference~[2]. It motivates the inequality by its role in diameter bounds and product theorems, and says that a relaxed orbit--stabilizer theorem yields a lower bound for \(|AA^{-1}\cap C_H(g)|\). No precise lower bound is supplied. The passage also suggests that buildings may help improve growth and diameter bounds; this is contextual speculation, not an additional formal assertion.
\subsection{Short English statement}
The problem asks for bounds on how many elements of a generating set lie in a regular semisimple conjugacy class, with the constants depending as mildly as possible on the group rank, preferably polynomially. It also asks whether character theory can improve these bounds. The group hypotheses, rank convention, and a damaged conjugacy-class alternative remain unresolved. The supplied record labels the problem open; this draft does not independently verify that status.
\subsection{Sources}
\begin{itemize}
\item Supplied local UnsolvedMath registry record, permanent ID \(30005880\), title: \emph{Conjugacy-Class Intersections with Generating Sets}; exact source problem code: \texttt{catalogue entry}. The supplied passage is headed ``Problem 2.'' The record belongs to the Oberwolfach Reports Open Problems collection and gives proposed year 2024; no individual proposer is supplied.
\item Supplied report citation: \emph{Mini-Workshop: Growth and Expansion in Groups} (2024), Oberwolfach Reports, pp.~1005--1040, DOI \texttt{10.4171/owr/2024/17}. \url{https://doi.org/10.4171/owr/2024/17}.
\item Supplied literature source, labelled ``OWR report'': \url{https://ems.press/content/serial-article-files/49475}.
\item Supplied literature source, labelled ``conjugacy-class estimates'': \url{https://arxiv.org/abs/2308.09197}. No authors, title, or precise supporting result are supplied in the record.
\item Literature assessment checked at \texttt{2026-08-22}: the supplied curation labels the problem open and states that the report records only a bound of order rank to a polynomial in rank squared, while polynomial rank dependence remains open. This is supplied dated curation derived from the OWR material, not an independently verified literature assessment. The record expressly requests independent verification of time-sensitive status and cited claims.
\item The report passage cites reference~[2] for bounds involving arbitrary varieties, but its bibliographic entry is absent from the supplied record. Its identity cannot be inferred from the other supplied links.
\item The registry labels its statement \texttt{corrected\_verified} and says that only formula delimiters, notation, grammar, and typography were repaired, with content and scope preserved. This is a supplied editorial qualification, not independent mathematical verification.
\end{itemize}
\end{document}
